% !TEX program = xelatex
% Research design manuscript. No experimental results are claimed.
\documentclass{misisarticle}

\grnti{28.23.25}
\asjc{1702, 2613}
\oecd{1.2}
\title{Прогнозирование оттока пользователей цифровой подписки: сравнение логистической регрессии и анализа выживаемости}
\misisauthor{Роман Алексеевич Даютов}{кафедра инженерной кибернетики НИТУ «МИСИС»}{Москва, Россия}{m2312161@edu.misis.ru}

\abstract{Рассматривается прогнозирование отключения автопродления цифровой подписки. Предлагается сравнить логистическую регрессию и модель Кокса на обезличенных исторических данных подписочного сервиса. Сравнение планируется провести на горизонтах 7, 14 и 30 дней с временным разделением выборок, оценкой качества ранжирования и калибровки вероятностей. Практическая цель исследования заключается в выборе подхода для адресных удерживающих коммуникаций.}
\keywords{прогнозирование оттока, цифровая подписка, логистическая регрессия, анализ выживаемости, модель Кокса, правое цензурирование}

\begin{document}
\maketitle

\section{Введение}

Для подписочного цифрового сервиса своевременное выявление риска отключения позволяет заранее сформировать аудиторию удерживающих коммуникаций. Пользователю с повышенным риском может быть направлено предложение со скидкой посредством SMS или мобильного приложения. Для выбора момента коммуникации существенны как вероятность отказа, так и горизонт, на котором ожидается событие.

В настоящем исследовании отток определяется как отключение пользователем автопродления основной подписки. Событие необходимо отличать от окончания оплаченного доступа: после отказа от продления пользователь некоторое время может продолжать пользоваться сервисом. Выбор или отключение отдельного онлайн-кинотеатра также не является целевым событием.

Распространённая постановка прогнозирования оттока основана на бинарной классификации: требуется оценить вероятность отключения в течение заданного числа дней. Альтернативная постановка использует анализ выживаемости и учитывает время до события, а также незавершённые наблюдения. Цель работы заключается в разработке протокола сравнения этих подходов на одинаковых исторических срезах и в оценке их пригодности для раннего выявления риска отключения.

В статье представлены постановка задачи, обзор связанных работ, состав данных и методика запланированного эксперимента. Полученные численные результаты и выводы о превосходстве конкретного подхода на данном этапе не заявляются.

\section{Анализ существующих работ}

В исследовании Burez и Van den Poel~\cite{burez2007} рассматривается прогнозирование оттока клиентов платного телевидения и применение аналитических моделей в адресных маркетинговых кампаниях. Работа близка к рассматриваемой предметной области: в обоих случаях анализируются подписочная услуга и возможность использования прогноза для удержания. В числе рассматриваемых методов используется логистическая регрессия.

Kimitei и соавторы~\cite{kimitei2025} исследуют анализ выживаемости для прогнозирования оттока клиентов коммунального сервиса. Авторы рассматривают модель пропорциональных рисков Кокса и аддитивную модель Аалена. Для настоящей работы существенны учёт цензурирования, оценка риска во времени и интерпретация факторов, связанных с отключением. Переносимость результатов на цифровую подписку требует самостоятельной проверки.

В препринте Khan~\cite{khan2020} рассматривается отток пользователей World of Warcraft, распространяемой по подписочной модели. В работе сочетаются анализ выживаемости с оценкой Каплана--Майера и бинарная классификация, включая логистическую регрессию. При этом исследование не является прямым сравнением индивидуального прогноза модели Кокса с логистической регрессией на одинаковых горизонтах.

В работах студентов кафедры представлены смежные прикладные задачи. Личко~\cite{lichko2023} сравнивает модели прогнозирования по историческим временным данным и оценивает прикладную стратегию. Соседка и Ибрагимов~\cite{sosedka2024} исследуют анализ цифровой аудитории и формирование рекомендаций. Эти публикации используются как примеры организации прикладного исследования; задачи оттока и времени до отключения в них не решаются.

Таким образом, предметом настоящего исследования является сопоставление фиксированных горизонтов классификации с вероятностями, полученными из survival-модели, при едином определении события, составе признаков и временном протоколе проверки. Исследование не предполагает, что один из подходов заранее обладает преимуществом.

\section{Постановка задачи и данные}

\subsection{Единица наблюдения и признаки}

Планируется использовать обезличенные корпоративные данные подписочного сервиса. История отключений автопродления доступна с июля 2025 года. Фактические границы выгрузки, число пользователей и доля событий будут зафиксированы после проверки полноты источников.

Исходная витрина содержит ежедневные состояния подписок. Дата среза $t$ соответствует концу календарного дня. Прогноз строится для подписок, у которых на этот момент автопродление включено и целевое отключение ещё не произошло. Признаки должны отражать только сведения, доступные к моменту формирования прогноза.

Предполагаются следующие группы признаков: время с подключения основной подписки; тариф, стоимость, период оплаты и скидка; выбранный онлайн-кинотеатр и время с момента его выбора; последние доступные месячные показатели ARPU и DOU. Дополнительно рассматриваются возраст, пол и доступные показатели просмотров.

Показатели активности планируется получать из агрегатов просмотров START. Возможные признаки включают давность последнего просмотра, число дней с просмотром и суммарные минуты за предыдущие 7 и 30 дней. При неподтверждённой полноте источника такие значения должны оставаться пропущенными. Отсутствие данных не приравнивается к отсутствию просмотров. Аналогично отсутствие истории по другим кинотеатрам не означает нулевую активность пользователя.

Для месячных и нерегулярно обновляемых источников требуется учитывать фактическую доступность записи. Отчётный месяц показателя сам по себе не подтверждает, что показатель уже был известен в дату прогноза. Если журнал загрузок отсутствует, для эксперимента задаётся и документируется консервативное допущение о задержке публикации; отдельно оценивается чувствительность результатов к его изменению.

Идентификаторы клиента и подключения используются для соединения источников и контроля выборок. Номер телефона и технические идентификаторы не включаются в признаки модели и не публикуются в репозитории.

\subsection{Разметка события и цензурирования}

Для каждого исторического среза устанавливается разрешённый конец наблюдения $C_i$. Он ограничен достоверно доступной историей и датой, до которой разрешено использовать исходы в соответствующем историческом запуске. Дата отключения $T_i$ используется только для разметки, а не как входной признак.

Целевые значения survival-модели задаются парой:
\begin{equation}
 d_i=\min(T_i,C_i)-t_i,\qquad
 \delta_i=\mathbf{1}\{t_i<T_i\leq C_i\}.
\end{equation}
Если отключение не наблюдалось, $d_i=C_i-t_i$, а $\delta_i=0$. В витрине эти атрибуты обозначаются как \texttt{duration\_days} и \texttt{event\_observed}. В первом случае наблюдается точное время до события, во втором известно только, что отключения не было в течение наблюдавшегося периода. Это соответствует правому цензурированию.

Например, для среза 1 сентября отключение 11 сентября даёт длительность 10 дней и индикатор события 1. Если история доступна только до 21 сентября и отключения нет, длительность составляет 20 дней, а индикатор равен 0. Последнее наблюдение не означает, что пользователь сохранит автопродление навсегда.

Для логистической регрессии целевая переменная на горизонте $H$ равна 1 при отключении в интервале $(t_i,t_i+H]$ и 0 при подтверждённом отсутствии отключения на всём интервале. При недостаточном наблюдении исход остаётся неизвестным, если событие ещё не произошло. Поэтому \texttt{event\_observed} нельзя непосредственно использовать как целевую переменную классификации. Например, отключение через 45 дней означает отрицательный класс на горизонте 30 дней, хотя индикатор события survival-модели равен 1.

\section{Методология: сравнение двух подходов}

\subsection{Логистическая регрессия}

Логистическая регрессия оценивает вероятность отключения в течение фиксированного горизонта по признакам клиента на дату среза. Для горизонтов 7, 14 и 30 дней планируется обучить отдельные модели с регуляризацией. Параметры регуляризации подбираются на временной валидационной выборке.

Числовые признаки обрабатываются с учётом пропусков и при необходимости масштабируются; категориальные кодируются. Все параметры преобразований оцениваются только на обучающей выборке. Сильной стороной подхода является простота получения и интерпретации прогноза на конкретный горизонт. Ограничение заключается в необходимости отдельной разметки для каждого горизонта и в невозможности трактовать незавершённое наблюдение как обычный отрицательный класс.

\subsection{Модель Кокса}

В качестве основной survival-модели рассматривается регуляризованная модель пропорциональных рисков Кокса. Обучение использует признаки среза, длительность наблюдения и индикатор события. При оценивании учитываются как наблюдавшиеся отключения, так и интервалы, завершившиеся без зарегистрированного события.

Результат модели представляется функцией $S(H\mid x)$: вероятностью того, что после даты среза автопродление не будет отключено в течение $H$ дней при заданных признаках $x$. Вероятность отключения на этом горизонте получается как $1-S(H\mid x)$. Таким образом, из одной модели можно получить прогнозы на те же горизонты, что и у классификатора. Результат не следует интерпретировать как гарантированную дату будущего отключения.

Перед интерпретацией коэффициентов требуется проверить допущение пропорциональности рисков~\cite{kimitei2025}. При существенном нарушении рассматриваются изменение спецификации или альтернативная survival-модель. Возможность получать несколько горизонтов из одной модели сама по себе не гарантирует более высокое качество прогнозов.

\subsection{Протокол исторической проверки}

Проверка планируется с последовательным разделением на обучение, валидацию и тест. Для каждого исторического запуска обучающие исходы ограничиваются датой запуска: сведения о последующих отключениях не должны попадать в обучение. Финальная тестовая выборка формируется на более поздней дате и не используется для выбора признаков или параметров.

В начальном эксперименте для каждой выборки выбирается один опорный исторический срез. Внутри такой выборки каждый период подписки представлен не более одной строкой. Это позволяет использовать ежедневную витрину для подготовки признаков, не превращая ежедневные состояния одного клиента в набор независимых объектов обучения. Устойчивость результатов затем проверяется на дополнительных последовательных опорных датах.

Случайное перемешивание всех ежедневных строк не используется. Клиенты, продолжающие пользоваться сервисом, могут встречаться в разных временных выборках, что соответствует сценарию регулярного прогнозирования существующей базы. Однако будущие исходы, перекрывающие дату следующего запуска, ограничиваются соответствующим концом наблюдения. Для оценки неопределённости повторные наблюдения одного периода подписки учитываются совместно.

Основное сравнение проводится на одинаковых тестовых клиентах и горизонтах с достаточным периодом последующего наблюдения. Планируется оценивать ROC-AUC и PR-AUC, ошибку вероятностного прогноза Brier score, калибровку, а также precision и recall среди фиксированной доли клиентов с наибольшим прогнозируемым риском. Для survival-модели дополнительно может использоваться индекс согласованности; он не заменяет сравнение вероятностей на общих горизонтах.

Различия в метриках оцениваются с помощью повторных выборок на уровне периодов подписки с сохранением их зависимости между срезами. В отчёте фиксируются даты выборок, частота целевого события, доступность признаков и ограничения исходных данных. При использовании цензурированной тестовой выборки потребуются метрики, явно учитывающие цензурирование; обычные метрики классификации на неизвестных исходах не рассчитываются.

\subsection{Бизнес-применение}

Практический результат эксперимента должен показать, какой подход лучше выделяет аудиторию повышенного риска при заданном объёме коммуникаций и насколько его вероятности соответствуют наблюдаемой частоте отключений. Это позволяет обосновать дальнейший пилот удерживающей кампании.

Историческая проверка оценивает качество прогноза, но не доказывает, что отправка скидки снижает отток. Эффект предложения и его экономическая целесообразность должны проверяться отдельно в рандомизированном эксперименте с контрольной группой. На этапе настоящей статьи описывается методика проверки моделей; результаты такого бизнес-эксперимента не заявляются.

\begin{thebibliography}{9}

\bibitem{burez2007}
Burez~J., Van den Poel~D.
CRM at a pay-TV company: Using analytical models to reduce customer attrition by targeted marketing for subscription services.
\textit{Expert Systems with Applications}. 2007. Vol.~32, no.~2. P.~277--288.
DOI: \href{https://doi.org/10.1016/j.eswa.2005.11.037}{10.1016/j.eswa.2005.11.037}.

\bibitem{kimitei2025}
Kimitei~S., Agiro~D., Ni~S., Ni~H.
Predictability \& explainability of survival analysis in churn prediction.
\textit{Journal of Marketing Analytics}. 2025.
DOI: \href{https://doi.org/10.1057/s41270-025-00450-2}{10.1057/s41270-025-00450-2}.

\bibitem{khan2020}
Khan~S. Predicting Customer Churn in World of Warcraft.
Препринт arXiv:2006.15735. 2020.
\url{https://arxiv.org/abs/2006.15735}.

\bibitem{lichko2023}
Личко~Д.~А. Разработка стратегии торговли биткоином с использованием методов машинного обучения.
\textit{Искусственный интеллект в промышленных, коммерческих, медицинских и финансовых приложениях: сборник статей научно-технического семинара студентов кафедры инженерной кибернетики}. М.: НИТУ «МИСИС», 2023. С.~151--157.
\url{https://misis.ru/files/-/1a3a57deb59c39284cba56d0b54bb308/kik_2023.pdf}.

\bibitem{sosedka2024}
Соседка~А.~В., Ибрагимов~П.~И. Количественный и качественный анализ аудитории Telegram в разрезе рекомендаций с использованием больших языковых моделей.
\textit{Искусственный интеллект в промышленных, коммерческих, медицинских и финансовых приложениях: сборник статей научно-технического семинара студентов кафедры инженерной кибернетики}. М.: НИТУ «МИСИС», 2024. С.~116--120.
\url{https://misis.ru/files/-/7914c2bff86d1a41832cf9bc14e3de9b/kik_2024_1.pdf}.

\end{thebibliography}
\end{document}
